\chapter*{Resumen}

\section*{\tituloPortadaVal}

Este Trabajo de Fin de Grado presenta el diseño de NutriApp, una plataforma web pensada para centralizar el flujo de trabajo de los nutricionistas profesionales. La idea surgió al observar de cerca cómo trabaja un nutricionista en su día a día: su información estaba repartida entre cuadernos físicos, formularios de Google, Google Calendar, la plataforma de su clínica, notas del móvil y WhatsApp. Si se quería analizar el estado completo de un solo cliente, se tenía que cruzar datos de todas estas herramientas de forma manual; lo cual resulta ineficiente y es probable perder información por el camino.

A partir de este problema real, se propuso diseñar una aplicación web que reuniera todas estas funcionalidades en un único punto de acceso. La plataforma se estructura en dos paneles diferenciados: uno orientado al nutricionista, desde el que puede gestionar clientes, citas, planes nutricionales y notas privadas; y otro orientado al cliente, con una interfaz propia donde consultar su plan, registrar sus hábitos diarios, comunicarse con su nutricionista y hacer seguimiento de su progreso.

Además de sustituir las herramientas dispersas, NutriApp incorpora una funcionalidad innovadora: la creación de dietas personalizadas para cada cliente. Esta herramienta tan específica solo es posible de llevar a cabo con éxito cuando gran parte de la información relativa al usuario está conectada, siendo posible acceder a sus alérgenos, intolerancias, objetivos nutricionales, preferencias y hábitos entre otros. Además, al mismo tiempo, es de gran relevancia las normas, preferencias y restricciones que imponga su nutricionista al programa, ya que condicionarán al plan diseñado por la Inteligencia Artificial que construirá sobre los alimentos guardados y modelados sobre la base datos; ya que cada uno de ellos tiene etiquetas que los sitúan en ciertos grupos de alérgenos, familia de alimentos, y también especifican los macronutrientes y propiedades que contienen.

\section*{Palabras clave}

\noindent FIXEAR: Plataforma web, nutrición, gestión de clientes, base de datos relacional, alertas inteligentes, gamificación, adherencia, formularios personalizables, lista de la compra, reconocimiento de alimentos, marketplace, automatizaciones, centralización
